\scanpage{174}\mspage{119}
Une remarque similaire s'applique à $\underline{\unsure{\mathrm{Fibhot}|\mathrm{Fiblat}}}\,S$, où $W_S^d=W_S^f$ également. Le foncteur chgt de base induit
\[
\underline{\mathrm{Fibhot}}\,S\longrightarrow\underline{\mathrm{Fibhot}}\,S',\quad X\longmapsto X\times_S S'\,/S',
\]
compatible avec $\deleted{\unsure{W_S^d}\,\unsure{W_S^f}}$ les localiseurs relatifs pertinents (tous ceux envisagés jusqu'à présent coïncident sur ces catégories, avec les hyp. faites sur $W$), d'où
\[
f^*:(\underline{\mathrm{Fibhot}}\,S)W_S^{-1}\longrightarrow(\underline{\mathrm{Fibhot}}\,S')W_{S'}^{-1}.
\]

Les localisés cette fois ne sont pas \inserted{les} $\mathrm{HOT}_W$, mais les $\mathrm{HOT}^{\mathrm{lc}}_W$, donc on trouve une description sympathique de
\[
f^*:\mathrm{HOT}^{\mathrm{lc}}_W(S)\longrightarrow\mathrm{HOT}^{\mathrm{lc}}_W(S'),
\]
plus avantageuse \unsure{encore|même} qu'en \deleted{\unsure{fonct.}} termes de $\underline{\mathrm{Liss}}_0\,S$, $\underline{\mathrm{Liss}}_0\,S'$, car s'appliquant à des catégories \inserted{un peu} plus vastes (NB $\underline{\mathrm{Liss}}_0\subset\underline{\mathrm{Fibhot}}$).

Notons donc le

\uline{Corollaire 1} La formule de la prop.~1 \inserted{(p.~118)} reste valable si $X$ est supposé $\mathrm{Hot}$-fibrant sur $S$ (au lieu de lisse sur $S$). (On trouve alors des types d'homotopie relatifs \uline{loc. const.} sur $S,S'$.)

On a donc une description de $f^*$ qui s'applique à \unsure{toutes} les catégories du côté gauche
% archive sheet number: 173
